\chapter{Introduction to Imperative C}

\textbf{Readings:} CP:AMA 4.2--4.5, 5.2, 7.1, 10.2--10.4.

\section{Functional vs.\ Imperative Programming}

\begin{definition}
In the \emph{functional} paradigm, programs only use constant values that never change;
functions produce values and have no other effect. In the \emph{imperative} paradigm, a program
is a sequence of \emph{statements} (commands) that change the \emph{state} of the program.
\end{definition}

Racket is mostly functional (but has imperative features); C is primarily imperative.

\section{Side Effects}

\begin{definition}
An expression or function has a \emph{side effect} if it does more than produce a value: it also
changes the state of the program (or ``the world'').
\end{definition}

Printing is the simplest side effect: \code{printf} changes what is on the screen. In Racket,
\code{printf} produces \code{(void)}, which is useless as a value; we only call it for its effect.
To do several things in sequence in Racket we use \code{begin}, which evaluates each expression
in order and produces the value of the \emph{last} one (the values of the others are ignored):
\begin{lstlisting}[style=racket]
(define (noisy-sqr n)
  (begin
    (printf "squaring ~a\n" n)
    (* n n)))
\end{lstlisting}
Function bodies and \code{cond} answers in full Racket have an implicit \code{begin}.

\subsection{Statements in C}

C functions are written as sequences of statements:
\begin{itemize}
  \item an \emph{expression statement} is an expression followed by a semicolon; its value is
    discarded, so it is only useful for its side effects (e.g.\ \code{printf("hi\textbackslash n");});
  \item a \emph{compound statement} (block) is a sequence of statements in braces
    \code{\{ ... \}};
  \item \emph{control flow statements} (\code{return}, \code{if}, loops) change the order of
    execution (Unit~05).
\end{itemize}

\subsection{\code{void} functions}

A function that only has side effects and produces nothing has return type \code{void}. It may
use \code{return;} (no value) to exit early, or simply reach the end of its body.
\begin{lstlisting}[style=c]
// print_square(n) prints the square of n
// effects: produces output
void print_square(int n) {
  printf("%d squared is %d\n", n, n * n);
}
\end{lstlisting}

\begin{important}[Documenting side effects]
If a function has a side effect, its documentation must contain an \textbf{effects:} section,
e.g.\ ``effects: produces output'' or ``effects: modifies *p'' or ``effects: modifies the
global counter''.
\end{important}

\section{State}

The biggest difference between the imperative and functional paradigms is the existence of side
effects that change the state.

\begin{definition}
The \emph{state} of a program is the value of all of its data at a moment in time.
\end{definition}

In a functional program, the result of a function call depends only on its arguments. In an
imperative program it may also depend on the state, so the same call can give different results
at different times.

\section{Mutation}

\begin{definition}
\emph{Mutation} is changing the value of a variable (or of data in memory).
\end{definition}

In Racket, \code{set!} mutates a variable:
\begin{lstlisting}[style=racket]
(define count 0)
(define (increment!)
  (set! count (add1 count))
  count)
(increment!)  ; => 1
(increment!)  ; => 2   (same expression, different value!)
\end{lstlisting}

\subsection{Variables in C}

A definition \emph{without} \code{const} creates a mutable \emph{variable}:
\begin{lstlisting}[style=c]
int count = 0;           // mutable global variable
const int limit = 10;    // constant: cannot be changed

int increment(void) {
  count = count + 1;     // assignment
  return count;
}
\end{lstlisting}

\begin{important}[Assignment]
\code{x = e} is the \emph{assignment operator}. It evaluates \code{e}, stores the result in
\code{x}, and (as an expression) also produces that value. Initialization (\code{int x = 5;})
and assignment (\code{x = 5;}) look similar but are different: \code{const} variables may be
initialized but never assigned.
\end{important}

\begin{aside}
Because assignment is an expression, \code{x = y = 0;} works (assigns right to left), and
\code{if (x = 0)} compiles. The latter is almost always a bug. Some people write
\code{if (0 == x)} so that the typo \code{0 = x} gives a compile error.
\end{aside}

\subsection{Uninitialized variables}

In C, a variable can be defined without an initial value: \code{int x;}. A \emph{global}
uninitialized variable is automatically set to 0, but a \emph{local} one contains whatever
arbitrary bits were left in that memory. Reading it is a bug. Always initialize variables.

\subsection{Compound assignment and increment}

\begin{center}
\begin{tabular}{ll}
  \toprule
  Shorthand & Meaning \\
  \midrule
  \code{x += y} & \code{x = x + y} \\
  \code{x -= y} & \code{x = x - y} \\
  \code{x *= y}, \code{x /= y}, \code{x \%= y} & similarly \\
  \code{++x} & increment $x$, produce the \emph{new} value \\
  \code{x++} & increment $x$, produce the \emph{old} value \\
  \code{--x}, \code{x--} & likewise for decrement \\
  \bottomrule
\end{tabular}
\end{center}

\begin{lstlisting}[style=c]
int x = 5;
int a = x++;   // a is 5, x is 6
int b = ++x;   // b is 7, x is 7
\end{lstlisting}
Prefer \code{++x} when the value is not used. Never use \code{x++} twice on the same variable in
one expression (e.g.\ \code{x = x++ + x++;}): that is undefined behaviour.

\section{Global Mutable State: Why It Is Dangerous}

Mutable global variables make programs harder to reason about and test, because any function
anywhere could change them. They also increase coupling between modules. In CS 136 they are
allowed but discouraged; most state should be passed to functions explicitly (with pointers,
Unit~06). When a module does need private state, make the variable \code{static} so only that
module can change it:
\begin{lstlisting}[style=c]
// counter.c
static int count = 0;      // module-scope state, hidden from clients

// increment() increases the counter by one and produces the new value
// effects: modifies the module's counter
int increment(void) {
  ++count;
  return count;
}

void reset(void) {
  count = 0;
}
\end{lstlisting}

\section{Mutable Structures}

Fields of a non-\code{const} structure can be mutated with assignment:
\begin{lstlisting}[style=c]
struct posn p = {1, 2};
p.x = 5;           // p is now {5, 2}
p.y += p.x;        // p is now {5, 7}

struct posn q = p; // copies all fields; q is independent of p
q.x = 0;           // p.x is still 5
\end{lstlisting}
Assigning one structure to another copies every field. This matters a lot once we pass
structures to functions (Units~05 and~06).

\section{Testing Imperative Code}

With side effects, testing must check both the produced value and the changed state:
\begin{lstlisting}[style=c]
int main(void) {
  reset();
  assert(increment() == 1);
  assert(increment() == 2);
  reset();
  assert(increment() == 1);
}
\end{lstlisting}
Testing output requires comparing the printed output against an expected output file
(Unit~07).

\section{Summary}
\begin{itemize}
  \item Imperative programs are sequences of statements that change state.
  \item A side effect is any change of state (output, mutation); document it under
    \emph{effects}.
  \item Variables without \code{const} are mutable; \code{=} assigns; \code{+=}, \code{++x},
    \code{x++} are shorthands.
  \item Always initialize local variables; avoid mutable globals (or make them \code{static}).
  \item Structure assignment copies every field.
\end{itemize}
